\section*{Chapitre \ref{ch:s1:short-term-futures-strategies} --- Stratégies sur contrats à terme à court horizon}

\begin{solution}{exo:s1:short-term-futures-strategies:1}
$0.01/100 = 1$ point de base. Un aller-retour à une fourchette de 1.08 tick paie deux fois la moitié de la fourchette~: 1.08 points de base.
\end{solution}

\begin{solution}{exo:s1:short-term-futures-strategies:2}
$50 \times 0.3 = 15$ points de base. Huit événements par an multiplient par $\sqrt{8}$ le ratio de Sharpe propre à chaque événement~; avec l'espacement de 32 jours de séance du chapitre,
$\sqrt{252/32} = 2.81$.
\end{solution}

\begin{solution}{exo:s1:short-term-futures-strategies:3}
La dérive de tendance injectée est répartie uniformément sur la séance, si bien que plus la plage est longue, plus elle a déjà capté de la direction de la journée au moment où la cassure
se déclenche, et la cassure ne prend position qu'au plus une fois par jour, si bien que l'entrée plus tardive coûte peu.
\end{solution}

\begin{solution}{exo:s1:short-term-futures-strategies:4}
$1~\%/\sqrt{39} = 16$ points de base~; un vrai effet de 50 points de base donne $t = 0.5/0.16 = 3.1$ en moyenne (2.7 dans l'échantillon).
\end{solution}

\begin{solution}{exo:s1:short-term-futures-strategies:5}
Son edge brut par transaction, la part du mouvement qui rebondit, est inférieur à une demi-fourchette à l'entrée et à une demi-fourchette à la sortie~; avec plusieurs transactions par jour, les
coûts croissent avec le nombre de transactions, alors que l'edge, lui, non.
\end{solution}

\begin{solution}{exo:s1:short-term-futures-strategies:6}
La somme quotidienne de la séance est multipliée par environ $(1 - 0.1) = 0.9$, donc sa variance est multipliée par 0.81~; avec les trois quarts de la variance dans la séance, la
volatilité du jour est multipliée par $\sqrt{0.25 + 0.75 \times 0.81} = 0.926$.
\end{solution}

\begin{solution}{exo:s1:short-term-futures-strategies:7}
Avec un seuil de deux écarts-types, le fade rapporte 5.18 points de base par jour avant coûts sur 7.3 transactions par jour~: il atteint le seuil de rentabilité à 0.35 points de base
par côté, en deçà des 0.54 du tape. Une entrée passive à l'autre bout de la fourchette sur une jambe la ramènerait près du seuil~; les deux jambes passives, avec des exécutions
incertaines, c'est ce que fait un fournisseur de liquidité.
\end{solution}

\begin{solution}{exo:s1:short-term-futures-strategies:8}
Quarante événements ont une erreur type d'environ 16 points de base~: une moyenne de 40 est compatible avec une dérive réelle comprise entre environ 10 et 70, y compris une
dérive affaiblie au tiers de sa valeur. Cinq ans ne permettent pas de confirmer que l'edge est intact~; l'a priori de l'article, les coûts et le degré d'encombrement de la stratégie comptent davantage.
\end{solution}

\begin{solution}{pb:s1:short-term-futures-strategies:1}
\begin{enumerate}
\item Barres d'une minute sur dix ans~: volatilité en U, gap overnight, dérive de tendance quotidienne, rebond et \omterm{def:s1:short-term-futures-strategies:pre}{dérive pré-annonce} arrêtée après la cinquième année.
\item La moitié de la fourchette cotée pondérée par le temps, sur une séance du tape~: 0.54 points de base.
\item Une position sur le premier mouvement hors de la plage des premières minutes~; un fade sur les mouvements courts qui rebondissent.
\item La journée de tendance~; le rebond.
\item 0.35, 0.55 et 1.25 après coûts pour 15, 30 et 60 minutes.
\item 3.78, 2.86, 1.82 et 0.59 avant~; $-3.31$, $-1.50$, $-0.71$ et $-0.78$ après, pour des seuils de 1.5 à 3.
\item La cassure ne prend position qu'une fois par jour pour un gros edge~; le fade prend position souvent pour un petit.
\item Les fournisseurs de liquidité qui gagnent la fourchette.
\item Un mouvement avant une annonce programmée, quel que soit son contenu~; de forts rendements des actions américaines avant les décisions du FOMC, aucun sur les Treasuries ni avant d'autres
annonces.
\item Les huit réunions du FOMC de 2026.
\item 40.6 et 18.8 points de base par événement, ratios de Sharpe de 1.21 et 0.58.
\item Huit événements par an donnent une erreur type d'environ 16 points de base sur cinq ans.
\item \textbf{Résultat nommé.} La détention pré-annonce a rapporté 40.6 points de base par événement après coûts ($t = 2.7$) et un ratio de Sharpe de 1.21 avant
l'encombrement injecté, et 18.8 points de base ($t = 1.3$) et 0.58 après lui.
\item Rotation face à l'edge, taille de l'échantillon face à la décroissance, et le mécanisme.
\item Seulement les dates de réunion programmées, les heures de publication telles que prévues, les coûts, et assez d'années pour mesurer la décroissance.
\item Momentum post-annonce.
\item Le chapitre~13 a mesuré où tombent le rendement et le volume de la journée~; le présent chapitre prend position à l'intérieur de la séance sur ces régularités.
\item Exécutions passives et un coût inférieur à la taille du rebond, ce qui revient à être un fournisseur de liquidité.
\item Il a un signe quelle que soit la décision~; une information ferait bouger les cours dans le sens de la nouvelle.
\item Immédiateté et risque autour de moments prévisibles.
\end{enumerate}
\end{solution}

\begin{solution}{iq:s1:short-term-futures-strategies:1}
Elle prend position dans le sens du premier mouvement hors de la plage des premières minutes de la séance~; elle échoue les journées en range, quand les cassures s'inversent,
et lorsque beaucoup de traders placent leurs stops aux mêmes niveaux.
\end{solution}

\begin{solution}{iq:s1:short-term-futures-strategies:2}
Il faut la mesurer seulement après la date de publication, avec la même définition~; comparer sa taille à celle publiée et à son erreur type~; contrôler
d'autres marchés et d'autres annonces comme témoins~; et vérifier si le positionnement avant les réunions a augmenté.
\end{solution}

\begin{solution}{iq:s1:short-term-futures-strategies:3}
Coûts~: le nombre de transactions multiplié par la fourchette et les frais, face à l'edge brut par transaction~; ensuite, si les exécutions sont réalistes (position dans la file, sélection
adverse) et si les prix utilisés étaient négociables.
\end{solution}

\begin{solution}{iq:s1:short-term-futures-strategies:4}
Barres ou messages à horodatage exact, la fourchette à chaque transaction, un modèle d'exécution (fill model) pour les ordres passifs, les calendriers de séance et les rolls, les heures programmées des événements,
et la latence entre le signal et l'ordre.
\end{solution}

\begin{solution}{iq:s1:short-term-futures-strategies:5}
Un saut à l'annonce allant contre la position, des fourchettes plus larges et des carnets d'ordres plus minces autour, et une sortie encombrée si beaucoup de participants détiennent la même position.
\end{solution}

\begin{solution}{iq:s1:short-term-futures-strategies:6}
Détenir $-\operatorname{sign}(r_t)$ pendant la minute suivante rapporte $-\operatorname{sign}(r_t)(e_{t+1} + \theta e_t)$, avec pour espérance
\[
-\theta\,E[\operatorname{sign}(r_t)\,e_t] = -\theta\sigma\sqrt{2/\pi}\big/\sqrt{1 + \theta^2}
\]
pour des chocs gaussiens, positive pour $\theta < 0$. Avec un coût $k$ dans chaque sens, le seuil de rentabilité est atteint quand $2k$ égale cette espérance~: pour
$\theta = -0.1$, l'edge représente environ 8~\% de la volatilité d'une minute, qui doit dépasser une fourchette.
\end{solution}
